\documentclass[12pt]{article}
\usepackage{amsmath, amssymb, amsthm}
\begin{document}

\title{UET: Математическая формализация}
\author{UAF Economic Theory}
\date{Версия 1.0}
\maketitle

\section{Пространство состояний агента}
Агент $i$ в момент $t$ описывается вектором:
\[
\mathbf{T}_i(t) = \begin{pmatrix}
t_{self,i}(t) \\ t_{market,i}(t) \\ t_{obligation,i}(t) \\ t_{captured,i}(t)
\end{pmatrix},\qquad
\sum_{k} t_{k,i}(t) = 24 \text{ ч/сутки}.
\]

Глобальная консервация времени:
\[
T_{past}(t) + T_{now}(t) + T_{future}(t) = \text{const}.
\]

\section{Матрёшечная иерархия}
Уровни $l=0,\dots,L$. Коэффициент интеграции $\alpha_{i,l}\in[0,1]$.
Эффективность уровня $l$:
\[
E_l = \left( \prod_{i\in l} \alpha_{i,l} \right)^{1/N_l} \cdot \frac{\Delta T_{integrated,l}}{\Delta T_{cost,l}} .
\]

\section{Ключевые уравнения}
Богатство и власть:
\[
Wealth_i = \sum_j \alpha_{ij} \cdot t_{future,j}^{controlled},\qquad
P_i = \sum_j \alpha_{ij} (\Delta t_{captured,j} - \Delta t_{self,j}).
\]

Динамика распределения времени:
\[
\frac{d\mathbf{T}_i}{dt} = \mathbf{F}(\mathbf{T}_i,\boldsymbol{\alpha},\mathbf{E}) + \boldsymbol{\xi}_i(t).
\]

Захват времени:
\[
\frac{d t_{captured,i}}{dt} = \beta_i (1-\alpha_{self,i}) A_i - \gamma_i t_{self,i}.
\]

Интеграция в матрёшку:
\[
\frac{d\alpha_{i,l}}{dt} = \kappa (E_l - c_i) \alpha_{i,l}(1-\alpha_{i,l}) - \delta_{decay}(1-\alpha_{i,l}).
\]

Создание/разрушение будущего времени:
\[
\frac{d T_{future}}{dt} = r\,T_{past} - \pi M - \lambda T_{future}.
\]

\section{Макроэкономические агрегаты}
Net Free Time Index:
\[
\text{NFTI} = \frac{1}{N}\sum_i (t_{self,i} - \theta\, t_{captured,i}).
\]

Реальный экономический рост:
\[
g_{real} = \frac{d}{dt}\log\left( \sum_l E_l \cdot T_{past,l} \right).
\]

Инфляция:
\[
\pi = \max\left(0,\; \frac{\Delta M - \Delta T_{real}}{\Delta M}\right).
\]

\section{Устойчивость}
Стационарное состояние: $\frac{d\mathbf{T}}{dt}=0,\ \frac{d\boldsymbol{\alpha}}{dt}=0$.
Условие устойчивости матрёшки:
\[
E_l > 1 + \frac{\text{cost of integration}}{\text{benefit of scale}}.
\]

\section*{Заключение}
Модель может быть реализована как агент-ориентированная симуляция или система ОДУ. Параметры требуют эмпирической калибровки.

\end{document}
